% SPDX-FileCopyrightText: 2026 Will Cook
% SPDX-License-Identifier: CC-BY-4.0
\section{Review changes without changing their meaning}\label{sec:review}
\papersectiontarget{exposition-review}

A smoother sentence can make a different claim. Compare the proposal with
the current statement and proof before judging its style: what may vary,
which objects must be shared, and which direction of implication is justified?
The following cases show how a changed noun, verb or omitted condition can
alter those answers.

Record whether the proposal was accepted, revised, rejected or left pending
a named check, retaining the exact passages and the decision's source. A
text match locates a change; it does not establish equivalence. A missing
argument requires mathematical review, and acceptance of a local revision
does not adopt a universal writing rule.

\subsection{Replace a local term without losing its definition}

In the R7 \#249 revision, a private growth adjective was replaced by its
defining condition:
\begin{quote}
\emph{Before:} tempered integer carry orbit.\par
\emph{Selected revision:} integer carry sequence $u$ with
$u(N)/2^N\to0$.
\end{quote}
The gain is specific. A reader no longer has to recover a private definition
to know the required rate. It would be incorrect to replace this condition by
the more familiar phrase ``subexponential growth'', which says something
stronger. Other revisions retained standard terms such as upper Banach
density because those terms denote the properties the paper actually uses.
The test is definitional: substitute the proposed term's meaning back into
the claim, then check that the same sequences qualify.

Removing a shorthand can also remove the visible source of its assumptions.
In \#243, ``positive exact state'' bundled recurrence and positivity
conditions. A proposed replacement referred to the same assumptions after
that definition had disappeared. The integrating reviewer instead stated the
recurrences for the first implication, then explicitly added natural-number
domains and $a_n>1$, $C_n>0$ for the eventual consequence. In a second passage,
the reviewer removed ``earlier'' from ``earlier decreases'': the original
qualification did not restrict the time at which decreases could occur.

These R7 selections are not the original proposals or later formal-correspondence
and rendering checks. The replacement must carry the same assumptions and
quantifiers, not merely read more easily.

\subsection{Check the ordinary words that carry scope}

An unchanged formula can acquire a stronger claim from its surrounding prose.
The R11 \#251 revision replaces ``uses all late indices'' by
``permits changes at every index after the prefix'' in a comparison
construction~\cite[comparison construction]{paper251long}. The second describes allowed positions, not a promise that
every one is changed. In the sparse construction, the containing set $S$ and
target interval are fixed before the target is chosen; the correction depends
on that target, and its support may occupy only part of $S$
\cite[sparse-perturbation proposition]{paper251}.

The same revision replaces the claim that corrections can ``grow arbitrarily
slowly'' by an eventual upper bound chosen in advance. Given any prescribed
function tending to infinity, the corrections can eventually be bounded by
it. This does not say that the corrections themselves tend to infinity.
Read the choice order as an instruction: prescribe the bound first, then
construct corrections subject to it. This exposes what the short phrases
obscured. Expand the expression whose ordinary reading changes that order
or turns permission into a requirement.

\subsection{Rejecting an attractive universal rule}

Earlier reviews of the \#243 manuscript illustrate why historical advice
needs its original setting. One round asked for a cubic irrationality example
near the beginning. A subsequent round explicitly retained the live
bounded-negative result as the paper's lead and rejected an invariant rule
that an irrationality result must always come first. Later editorial routes
treated page targets as subordinate to the proof's dependencies.

The retained lesson is to choose a coherent principal contribution and explain
its relation to the named problem. It does not prescribe a theorem category
or a fixed page on which every proof must finish. A conditional reduction
can be the useful result; an unconditional statement can be too weak to carry
the paper. The choice requires comparing what the results actually say. The
history record preserves the successive recommendations instead of combining
them into contradictory commands.

This distinction also protects the longer research record. Shortening a
paper can be appropriate when the subordinate details have a complete,
checked destination. Regenerating the long record from the shortened version
would lose the very argument that made the shorter presentation possible.
The two documents must be compared in both directions: a stronger hypothesis,
a repaired endpoint or a newly explained limiting step may appear first in
either one.

\subsection{Distinguish clarification from a proof repair}

An earlier \#251 review found more than an unclear sentence. A condensed
sparse-construction proof used a buffer quantity without assigning it a
sequence coordinate or including it in the value, size and capacity accounting.
The original returned memorandum did include those coordinates and their
contribution to the weighted sum. The accepted ordinary repair restored the
assignments and the accounting across all indices.

The distinction matters for both mathematics and attribution. The condensed
proof needed repair; the recoverable original showed that the loss had occurred
during shortening or integration, not in that construction. This repair was
undertaken in its own mathematical review. An exposition-only pass should
identify such a defect and refer it, not silently supply the missing proof.
A later sharp companion construction was still unreviewed at its recorded
disposition. Sharing a bundle with the repaired proof did not give it the
same status.

Other cases expose similar changes of meaning in small amounts of prose.
In \#1049, language suggesting that a degree bound was inevitable was repaired
to state a sufficient inequality and preserve the qualification contributed
by a remaining factor. In \#1041, the revision retained the fixed-polynomial
scope instead of turning it into a freely varying family. For mixed
irrationality criteria in \#257, two unbounded sets of candidate indices
may be disjoint: the even and odd integers provide an elementary example.
The proof needs one index satisfying both conditions. The source adds the
two nonnegative, normalised errors on the same finite distribution. Their
sum has mean below one, so it is below one at some sampled index.
Nonnegativity puts both errors below one there; no independence is needed~\cite[common-index argument]{paper257}.
These are mathematical checks even when the edit is presented as compression or style.

Compare the Lean proposition with the prose, including definitions,
hypotheses and quantifiers. A checked ingredient need not cover the assembled
argument. An ordinary proof may lack complete formalisation. Keep named inputs
and pending comparisons visible. The systems paper's records preserve these
distinctions~\cite{systems}; describing them does not confer proof status.

A late-September prose pass removed remarks naming mathematical inputs beside
the dependent results; the integrating agent restored them. Details of running
a check could recede, but the dependencies could not. Calling both kinds of
sentence workflow commentary had erased a mathematical qualification.

\subsection{Read explanatory words as mathematical claims}\label{sec:words-as-claims}

Test a mathematical adjective by substituting its definition. A ``convex
combination'' requires nonnegative weights whose sum is one. Checking only
their signs leaves half the condition untested. In \#1041, Abel summation
expresses $f(t\zeta)$ as
\[
 f(t\zeta)=\sum_{j=0}^{n-1}(t^j-t^{j+1})S_j,
 \qquad S_j=\sum_{k=0}^{j}a_k\zeta^k,
\]
where $n\ge1$, $f(z)=\sum_{k=0}^n a_kz^k$ and $f(\zeta)=0$.
For $0\le t<1$, the nonnegative coefficients sum to $1-t^n$.
Thus it is $f(t\zeta)/(1-t^n)$ that the calculation directly presents as a
convex combination of the $S_j$~\cite[section on trinomials]{paper1041}.
At $t=1$ this quotient is undefined, and the proof instead uses $f(\zeta)=0$.
Naming the normalised quantity checks the definition and the endpoint in the
same sentence. A root-distance calculation or a statement about segments to
zeros must likewise retain the particular geometric object it controls.

A shorter proof can also be clearer when its reason is already available.
The \#243 arithmetic normal form is
\[
                 Q(n)=m\binom{n+2}{3}+c,
                 \qquad m\in\mathbb Z_{>0},
\]
and $Q(n)$ is known to be integral at every sufficiently large integer $n$.
At any such nonnegative $n$, the binomial coefficient is integral, so
$c=Q(n)-m\binom{n+2}{3}$ is integral
\cite[reduction and the constant term]{paper243}.
The proposed explanation uses exactly those established facts. Replacing them
by ``$Q$ is an integer polynomial'' would be unsafe: an integer-valued polynomial
can have nonintegral coefficients, as $n(n-1)/2$ shows. The useful simplification
names the available arithmetic structure.

In the \#257 finite average, $L,d,T$ are positive integers. The residues of
$L,2L,3L,\ldots$ modulo $d$ repeat with period $d/\gcd(L,d)$; the sample
uses only the first $T$ terms. The inequality $d\le LT$ does not assert
that one whole period occurs. For $L=3,d=5,T=2$, it holds while the period
has length five. Split the sample into complete periods, possibly none,
and one remainder. The kernel weights of Section~\ref{sec:formula-choice}
are nonnegative, so that remainder contributes no more than a whole period.
Allowing one extra period supplies the error term after division by $T$
\cite[finite means for shifted divisor tails]{paper257}.
The example refutes the claim of a complete sampled period, not the estimate.

\subsection{Correct the source without rewriting its history}

The R6 \#1049 return attached a reader-motivation specimen from Knuth's notes
to \S1, item~18. Inspection of the original places the passage in item~12,
on printed page~3 (PDF page~5). Item~18 concerns a different matter. The
public lesson record retains both the reported locator and the correction.

A search hit may be a table of contents, another occurrence or an adjacent
column. Inspect the original passage, then correct the locator without crediting
the earlier return with that correction. Useful prose can survive a citation
repair; its usefulness does not excuse the error.

The same care applies to the role of the source. A paper may supply a theorem,
historical attribution, established terminology or an example of composition.
Those uses call for different claims. A paragraph about proof pacing does not
validate the proof being paced. A literature specimen establishes that an
author made a particular choice, not that this choice is universally preferable.
